\clearpage
\VerdierPhysicalPageAnchor{0223}%
\begingroup\small\setlength{\parskip}{0pt}\setlength{\emergencystretch}{1.5em}
\setlength{\abovedisplayskip}{4pt plus 2pt minus 2pt}%
\setlength{\belowdisplayskip}{4pt plus 2pt minus 2pt}%

\VerdierSemanticAnchor{ch3d-3-3-8-cont}
使得 $m=m_nm_{n-1}\cdots m_2m_1$。这立即给出 $F$ 的唯一性。为证明存在性，
只须注意，由引理 (3.3.1)，当 $m$ 有第二个分解
$m=m'_nm'_{n-1}\cdots m'_1$ 时，存在交换图：

\VerdierSemanticAnchor{ch3d-3-3-8-comparison-diagram}
\begin{center}
\resizebox{\textwidth}{!}{%
$\begin{tikzcd}[ampersand replacement=\&,column sep=1.7em,row sep=3.0em]
{}
  \& X_1[p+1] \arrow[r,"m_2"]
  \& X_2[p+2] \arrow[r]
  \& \cdots \arrow[r]
  \& X_{n-1}[p+n-1] \arrow[dr,"m_n"]
  \& {} \\
X[p]
  \arrow[ur,"m_1"]
  \arrow[r,"{m''_1}"]
  \arrow[dr,"m'_1"']
  \& X''_1[p+1]
     \arrow[u,"{f_1[p+1]}"]
     \arrow[d,"{g_1[p+1]}"']
     \arrow[r,"{m''_2}"]
  \& X''_2[p+2]
     \arrow[u,"{f_2[p+2]}"]
     \arrow[d,"{g_2[p+2]}"']
     \arrow[r]
  \& \cdots \arrow[r]
  \& X''_{n-1}[p+n-1]
     \arrow[u,"{f_{n-1}[p+n-1]}"]
     \arrow[d,"{g_{n-1}[p+n-1]}"']
     \arrow[r,"{m''_n}"]
  \& Y[p+n] \\
{}
  \& X'_1[p+1] \arrow[r,"m'_2"']
  \& X'_2[p+2] \arrow[r]
  \& \cdots \arrow[r]
  \& X'_{n-1}[p+n-1] \arrow[ur,"m'_n"']
  \& {}
\end{tikzcd}\quad\mathpunct{,}$%
}
\end{center}

\VerdierSemanticAnchor{ch3d-3-3-8-conclusion}
其中 $f_i$ 与 $g_i$，$1\leq i\leq n-1$，都是 $\mathcal A$ 的态射。
\endgroup
